\begin{tikzpicture}[
  font=\scriptsize,
  execute at begin node={\hyphenpenalty=10000\exhyphenpenalty=10000},
  mulai/.style={draw, rounded corners, fill=praditagreen!15, text width=11.2cm, minimum height=0.65cm, align=center},
  gerbang/.style={draw, rounded corners, fill=orange!12, text width=3.3cm, minimum height=1.05cm, align=center},
  keluar/.style={draw, rounded corners, fill=blue!8, text width=3.3cm, minimum height=1.05cm, align=center},
  masalah/.style={draw, rounded corners, fill=gray!35, text width=3.3cm, minimum height=1.05cm, align=center},
  hasil/.style={draw, rounded corners, fill=green!10, text width=3.3cm, minimum height=1.05cm, align=center},
  jawab/.style={font=\tiny, inner sep=1pt, fill=white},
  panah/.style={-{Stealth[length=2mm]}, thick}
]
  % Tiga gerbang mengurangi kandidat sebelum skor, karena skor tidak boleh menggantikan kewajiban.
  \node[mulai] (mulai) at (4.0,1.05) {\textbf{Delapan kandidat} API-01 sampai API-08, ditambah kewajiban INT-01 dan INT-02};

  \node[gerbang] (g1) at (0.0,0) {1. Diwajibkan aturan?};
  \node[gerbang] (g2) at (4.0,0) {2. Ada pemilik dan bukti masalah?};
  \node[gerbang] (g3) at (8.0,0) {3. Ada alternatif lebih murah yang memadai?};

  \node[hasil] (w1) at (0.0,-1.75) {\textbf{Jalur kewajiban}, tanpa skor\\{\tiny\mdseries INT-01, INT-02: KI-01, KI-02}};
  \node[masalah] (w2) at (4.0,-1.75) {\textbf{Telusuri dulu}\\{\tiny\mdseries API-08 Resep Apotek}};
  \node[keluar] (w3) at (8.0,-1.75) {\textbf{Rapikan atau pakai}\\{\tiny\mdseries API-05, API-07}};

  \node[gerbang] (s) at (8.0,-3.5) {4. Beri skor: nilai, usaha, risiko\\{\tiny\mdseries API-01, 02, 03, 04, 06}};
  \node[gerbang] (g5) at (4.0,-3.5) {5. Prasyarat dan kapasitas terpenuhi?};
  \node[hasil] (p) at (0.0,-3.5) {\textbf{Portofolio terpilih}\\{\tiny\mdseries API-06, API-01, API-02}};
  \node[keluar] (a) at (4.0,-5.2) {\textbf{Antrean} dengan pemicu tinjauan\\{\tiny\mdseries API-03, API-04}};

  \draw[panah] (mulai.south -| g1.north) -- (g1.north);
  \draw[panah] (g1) -- node[jawab, above] {tidak} (g2);
  \draw[panah] (g2) -- node[jawab, above] {ya} (g3);
  \draw[panah] (g1) -- node[jawab, right] {ya} (w1);
  \draw[panah] (g2) -- node[jawab, right] {tidak} (w2);
  \draw[panah] (g3) -- node[jawab, right] {ya} (w3);

  % Panah tidak dari gerbang 3 memutari kotak keluaran agar alur ke skor tetap terbaca.
  \draw[panah] (g3.east) -- ++(0.28,0) |- node[jawab, pos=0.25, right] {tidak} (s.east);
  \draw[panah] (s) -- (g5);
  \draw[panah] (g5) -- node[jawab, above] {ya} (p);
  \draw[panah] (g5) -- node[jawab, right] {tidak} (a);
\end{tikzpicture}
